%% kapa: exam=1 round=25 year=2014 period=1 form=A subject=economics subject_ko=경제원론 seq=23 number=63
%% source: Q-Net 감정평가사 기출문제 (cst00309) artlSeq=5101207 file=extracted/1008/5101207/2014년 제25회 감정평가사 1차시험 기출문제/1교시(A형).hwp
%% conversion: kordoc 4.24.51 -> tools/gen1.py (auto); review=proofread-sonnet; flags=none
\begin{question}{63}
중앙은행이 요구불예금에 대한 법정지급준비율을 인상하면 지급준비금(ㄱ), 요구불예금(ㄴ)과 통화량(ㄷ)은 각각 어떻게 변화하는가? (단, 현재 법정지급준비율은 10\%이고, 민간은행들은 초과지급준비금을 보유하지 않으며 가계는 현금을 보유하지 않는다.)
\begin{choices}
  \item ㄱ: 감소, ㄴ: 감소, ㄷ: 불변
  \item ㄱ: 감소, ㄴ: 증가, ㄷ: 감소
  \item ㄱ: 불변, ㄴ: 감소, ㄷ: 증가
  \item ㄱ: 증가, ㄴ: 불변, ㄷ: 감소
  \item ㄱ: 증가, ㄴ: 불변, ㄷ: 증가
\end{choices}
\end{question}
\answerall
\answernote{정답표에 1,2,3,4,5 모두 정답으로 표기(전원 정답)}
